\begin{derivation}[从\eqnrefs{eq:qcd-bare-coupling-msbar-dp68}到\eqnrefs{eq:qcd-alpha-lambda-one-loop-dp10}]
重整化群方程不是额外加入 QCD 的动力学定律，而是同一个裸理论在改变参数定义尺度 \(\mu\) 时必须保持不变的条件。关键点是：在维数正规化的 \(d=4-2\epsilon\) 中，耦合本身先带有工程维数；这个 \(-\epsilon g_s\) 项与 \(Z_g\) 的紫外极点相乘，才留下四维中的有限运行。

从正文
\begin{equation*}
g_{s,B}=\mu^\epsilon Z_g(g_s,\epsilon)g_s,
\qquad
Z_g
=1-\frac{\beta_0}{2}
\frac{g_s^2}{(4\pi)^2}
\frac1{\bar\epsilon}
+\order(g_s^4)
\end{equation*}
出发。裸参数 \(g_{s,B}\) 不依赖人为的重整化尺度，因此
\begin{equation*}
0=\mu\frac{\dd}{\dd\mu}\ln g_{s,B}
=\epsilon
+\frac{\beta_g^{(d)}}{g_s}
+\beta_g^{(d)}\frac{\partial}{\partial g_s}\ln Z_g,
\end{equation*}
其中
\begin{equation*}
\beta_g^{(d)}(g_s)
\equiv\mu\frac{\dd g_s}{\dd\mu}
=-\epsilon g_s+\beta_g^{(4)}(g_s).
\end{equation*}
这里 \(-\epsilon g_s\) 是 \(d\) 维中的工程尺度项，而 \(\beta_g^{(4)}\) 是取四维极限后真正的量子运行。

只保留决定一圈有限项所需的最低阶。把 $\overline{\mathrm{MS}}$ 极点写成
\begin{equation*}
 \frac1{\bar\epsilon}
 \equiv
 \frac1\epsilon-\gamma_E+\ln(4\pi),
\end{equation*}
于是
\begin{align*}
\ln Z_g
&=-\frac{\beta_0}{2}
\frac{g_s^2}{(4\pi)^2}
\left[\frac1\epsilon-\gamma_E+\ln(4\pi)\right]
+\order(g_s^4),\\
\frac{\partial}{\partial g_s}\ln Z_g
&=-\frac{\beta_0 g_s}{(4\pi)^2}
\left[\frac1\epsilon-\gamma_E+\ln(4\pi)\right]
+\order(g_s^3).
\end{align*}
把 $\beta_g^{(d)}=-\epsilon g_s+\beta_g^{(4)}$ 代回裸尺度不变性。显式工程维数先给
\begin{equation*}
 \epsilon+\frac{-\epsilon g_s}{g_s}=0.
\end{equation*}
在剩下的乘积中，$\beta_g^{(4)}=\order(g_s^3)$ 乘上 $\partial_{g_s}\ln Z_g=\order(g_s/\epsilon)$ 只从两圈开始影响一圈方程，因此这一圈只需保留 $-\epsilon g_s$：
\begin{align*}
0
={}&\frac{\beta_g^{(4)}}{g_s}
+(-\epsilon g_s)
\left\{-\frac{\beta_0 g_s}{(4\pi)^2}
\left[\frac1\epsilon-\gamma_E+\ln(4\pi)\right]\right\}
+\order(g_s^4)+\order(\epsilon g_s^2)\\
={}&\frac{\beta_g^{(4)}}{g_s}
+\frac{\beta_0g_s^2}{(4\pi)^2}
\left[1-\epsilon\bigl(\gamma_E-\ln4\pi\bigr)\right]
+\order(g_s^4)+\order(\epsilon g_s^2).
\end{align*}
现在再取四维极限 $\epsilon\to0$，有限常数乘以 $\epsilon$ 的部分消失，而极点残数保留下来：
\begin{equation*}
0=\frac{\beta_g^{(4)}}{g_s}
+\frac{\beta_0g_s^2}{(4\pi)^2}
+\order(g_s^4).
\end{equation*}
因此
\begin{equation*}
\boxed{
\beta_g^{(4)}(g_s)
=-\frac{\beta_0}{16\pi^2}g_s^3
+\order(g_s^5).}
\end{equation*}
极点项 $1/\bar\epsilon$ 不能按普通代数量直接“约去”；有限的一圈运行来自严格极限
\begin{equation*}
 \lim_{\epsilon\to0}
 \epsilon\left(\frac1\epsilon-\gamma_E+\ln4\pi\right)=1.
\end{equation*}
这正是维数正规化中工程尺度项与紫外极点结合后留下有限反常尺度依赖的数学机制。

再用 \(\alpha_s=g_s^2/(4\pi)\)，链式法则给
\begin{align*}
\mu\frac{\dd\alpha_s}{\dd\mu}
&=\frac{g_s}{2\pi}
\mu\frac{\dd g_s}{\dd\mu}\\
&=-\frac{\beta_0}{2\pi}\alpha_s^2
+\order(\alpha_s^3),
\end{align*}
即正文\eqnrefs{eq:qcd-beta-alpha-one-loop-dp10}。

在活跃味数固定、\(\beta_0\) 为常数的尺度区间内，一圈方程可分离变量：
\begin{equation*}
\frac{\dd\alpha_s}{\alpha_s^2}
=-\frac{\beta_0}{2\pi}\dd\ln\mu.
\end{equation*}
从参考尺度 \(\mu_{\rm R0}\) 积到 \(\mathcal Q\)，
\begin{equation*}
\frac1{\alpha_s(\mathcal Q)}
-\frac1{\alpha_s(\mu_{\rm R0})}
=\frac{\beta_0}{2\pi}
\ln\frac{\mathcal Q}{\mu_{\rm R0}}.
\end{equation*}
把积分常数改写成具有动量量纲的 \(\Lambda_{\rm QCD}\)，定义
\begin{equation*}
\frac1{\alpha_s(\mathcal Q)}
\equiv
\frac{\beta_0}{4\pi}
\ln\frac{\mathcal Q^2}{\Lambda_{\rm QCD}^2},
\end{equation*}
便得到
\begin{equation*}
\boxed{
\alpha_s(\mathcal Q)
=\frac{4\pi}{\beta_0
\ln[\mathcal Q^2/\Lambda_{\rm QCD}^2]}.}
\end{equation*}
当 \(\beta_0>0\) 时，\(\mathcal Q\) 增大使分母增大，所以 \(\alpha_s\) 下降；这就是“渐近自由”在一圈阶的定量含义。反过来，\(\mathcal Q\) 接近 \(\Lambda_{\rm QCD}\) 时分母不再大，\(\alpha_s\) 失去小参数，说明微扰展开在红外失控；这一步本身并不等价于禁闭证明。
\end{derivation}
